\section{Равноускоренное движение}
\subsection{Неравномерное движение. Графики}

%%%%%%%%%%%%%%%%%%%%%%%%%%%%%%%%%%%%%%%%%%%%%%%%%%%%
%%%%%%%%%%%%%%%%%%%%%%%%%%%%%%%%%%%%%%%%%%%%%%%%%%%%
 
 \z[{\begin{wrapfigure}[5]{r}{0.3\textwidth}%
		\vspace{-30pt}%
		\begin{flushright}%
 			\includegraphics[width=0.3\textwidth]{2.1.1}%
		\end{flushright}%
	\end{wrapfigure}%
	}]%
	{На рисунке представлен график зависимости скорости тела от времени. Начальная координата тела $x(0) = 0$. Постройте графики зависимости ускорения и координаты тела, а также пройденного им пути от времени. Тело движется вдоль оси $Ox$.}
%
{}

%%%%%%%%%%%%%%%%%%%%%%%%%%%%%%%%%%%%%%%%%%%%%%%%%%%%

 \z[{\begin{wrapfigure}[6]{r}{0.3\textwidth}%
		\vspace{-10pt}%
		\begin{flushright}%
 			\includegraphics[width=0.3\textwidth]{2.1.2}%
		\end{flushright}%
	\end{wrapfigure}%
	}]%
	{На рисунке представлен график зависимости скорости тела от времени. Начальная координата тела $x(0) = 0$. Постройте графики зависимости ускорения и координаты тела, а также пройденного им пути от времени. Определите среднюю и среднепутевую скорости за первые $2.0$ и $5.0\, \text{с}$ движения. Тело движется вдоль оси $Ox$.}
%
{$v_\text{ср.2}=0.5\,\nicefrac{\text{м}}{\text{с}} $;~$v_\text{ср.п.2}=1.3\,\nicefrac{\text{м}}{\text{с}} $;~$v_\text{ср.5}=-0.9\,\nicefrac{\text{м}}{\text{с}} $;~$v_\text{ср.п.5}=1.7\,\nicefrac{\text{м}}{\text{с}}$}
	
%%%%%%%%%%%%%%%%%%%%%%%%%%%%%%%%%%%%%%%%%%%%%%%%%%%%
\z[{\begin{wrapfigure}[6]{r}{0.3\textwidth}%
		\vspace{-20pt}%
		\begin{flushright}%
 			\includegraphics[width=0.3\textwidth]{2.1.3}%
		\end{flushright}%
	\end{wrapfigure}%
	}]%
	{Материальная точка движется вдоль прямой. Постройте графики зависимостей скорости и координаты точки от времени, если график зависимости её скорости $v$ от координаты $x$ представляет собой:%	
	
	\lpu прямоугольник\lpuend%
	\lpu окружность (при определённом выборе масштабов осей).\lpuend%
	}
%
{}

%%%%%%%%%%%%%%%%%%%%%%%%%%%%%%%%%%%%%%%%%%%%%%%%%%%%	
\z[{\begin{wrapfigure}[1]{r}{0.4\textwidth}%
		\vspace{-55pt}%
		\begin{flushright}%
 			\includegraphics[width=0.3\textwidth]{2.1.4}%
		\end{flushright}%
	\end{wrapfigure}%
	}]%
	{По графику $a_x (t)$ (см. рисунок) постройте график $v_x(t)$, считая $v(0) = 0$.}
%
{}

%%%%%%%%%%%%%%%%%%%%%%%%%%%%%%%%%%%%%%%%%%%%%%%%%%%%
 \z{Автомобиль проехал по пятикилометровому участку дороги. Специальный прибор при этом записывал показания спидометра через каждые $10\, \text{метров}$. В результате получилась зависимость скорости автомобиля $v$ от пройденного пути $x$, показанная на рисунке. Оцените, за какое время $t$ автомобиль проехал эти пять километров.
		\begin{center}%
 			\includegraphics[width=0.41\textwidth]{2.1.5}%
		\end{center}%
}
%
{$t \approx 120 \pm 10 \,\text{с}$}
